\input{\string~/Documents/School/"UC Irvine"/ta_template2.0.tex}
\rh{2B}{7.2}
\chead{\today}
\graphicspath{ {\string~/Desktop} }
\usetikzlibrary{arrows}
%============ANSWER KEY TOGGLE===========
\newcommand{\ans}[1]{ \noindent {\color{red} \textbf{Answer:} #1}}
%\renewcommand{\ans}[1]{\vspace{3in}}
%==========================================
\begin{document}
\begin{center}
\textbf{Trigonometric Integration}
\end{center}

\begin{aun}[Strategies for Evaluating $\int \sin^m (x) \cos^n(x)dx$] For integrals of this type, we break into three cases:
\begin{description}
\item[Case 1, $n$ is odd:] Since $n$ is odd, we can write $n = 2k+1$ for some $k$. Then we substitute:
\[
\int \sin^m(x)\cos^{2k+1}(x)dx = \int \sin^m(x)(\cos^2(x))^{k} \cos(x)dx = \int \sin^m(x)(1-\sin^2(x))^{k} \cos(x)dx
\]
and continue by $u$-substitution on $\sin(x)$. 
\item[Case 2, $m$ is odd:] Since $m$ is odd, we can write $m = 2k+1$ for some $k$. Then we substitute:
\[
\int \cos^n(x)\sin^{2k+1}(x)dx = \int \cos^n(x)(\sin^2(x))^{k} \sin(x)dx = \int \cos^n(x)(1-\cos^2(x))^{k} \sin(x)dx
\]
and continue by $u$-substitution on $\cos(x)$.
\item[Case 3, $m,n$ is even:] Utilize the power reducing formulas:
\[
\sin^2(x) = \frac12 (1-\cos (2x)), \quad \cos^2(x) = \frac12 (1 + \cos(2x)), \mte{ and } \sin(x) \cos(x) = \frac12 \sin(2x)
\]
\end{description}
\end{aun}
\begin{aun}[Strategies for Evaluating $\int \tan^m (x) \sec^n(x)dx$] In a similar way, we'll break into two cases again:
\begin{description}
\item[Case 1: $n$ is even:] We can write $n = 2k$ for some $k$, save a factor of $\sec^2(x) = (1+\tan^2(x))$ to use the identity. 
\[
\int \tan^m(x) \sec^{2k} (x) dx = \int \tan^m(x) (\sec^2(x))^{k-1} \sec^2(x) dx = \int \tan^m(x) (1+\tan^2(x))^{k-1} \sec^2(x) dx
\]
and continue by $u$-substitution.
\item[Case 2: $m$ is odd:] Create a factor of $\sec(x)\tan(x)$ and use the identity as before; since $m$ is odd, there is $k$ for which $m = 2k+1$: 
\[
\int \tan^{2k+1}(x) \sec^n(x) dx = \int \tan^{2k}(x) \sec^{n-1}(x) \sec(x)\tan(x) dx = \int (\sec^{2}(x)-1)^k \sec^{n-1}(x) \sec(x)\tan(x) dx
\]
and proceed by $u$-substitution.
\end{description}
We also can use the following:
\[
\int \tan(x)dx = \ln \lrabs{\sec(x)} + C \mte{ and } \int \sec(x)dx = \ln \lrabs{\sec(x) + \tan(x)} +C
\]
\end{aun}



\begin{exx}
Compute the following integrals:
\[
\int \sin^2 (x)\cos^3 (x)dx \mte{ and } \int \tan^5 (x) dx
\]
\end{exx}
\ans{
This first integral is in the case where $n$ is odd and $n$ is even, so we can write:
\begin{alignat*}{2}
	\int \sin^2 (x)\cos^3 (x)dx &= \int \sin^2 (x) \cos^2(x) \cos(x) dx\\
						&= \int \sin^2 (x) (1- \sin^2 (x)) \cos(x) dx \jq{ By $\sin^2(x) + \cos^2(x) = 1$}\\
						&= \int u^2 (1 - u^2) du \jq{ Let $u = \sin(x)$, $du = \cos(x)dx$}\\
						&= \frac13 u^3 - \frac15 u^5 + C\\
						&= \frac13 \sin^3(x) - \frac15 \sin^5(x) + C
\end{alignat*}
For the latter integral, we just have to remember:
\begin{alignat*}{2}
\int \tan^5 (x) dx &= \int \tan(x) (tan^2(x))^2 dx \\
			&= \int \tan(x) (\sec^2 (x) - 1)^2 dx \\
			&= \int \tan(x) \lrp{\sec^4(x) -2 \sec^2(x) + 1}dx\\
			&= \int \sec^4(x)\tan(x)dx - 2\int \tan(x) \sec^2 (x) dx + \int \tan(x)dx
\end{alignat*}
Each of these integrals can be solved using a different trick, so we can split into three cases to address each one:
\[
\int \tan(x) dx =  \ln \lrabs{\sec(x)} + C
\]
from above. The next, 
\[
\int \underbrace{\tan(x)}_{u} \underbrace{\sec^2 (x) dx}_{du} = \frac{u^2}{2} + C = \frac{\tan^2(x)}{2} + C
\]
admits a $u$ substitution. Finally, to conquer the beast, we have:
\[
\int \sec^4 (x) \tan(x) dx = \int \lrp{\underbrace{\sec(x)}_{u}}^3 \underbrace{ \sec(x) \tan(x) dx}_{du} = \frac{u^4}{4} + C = \frac{\sec^4(x)}{4} + C
\]
Putting the pieces together, we have:
\[
\int \tan^5(x) dx = \frac{\sec^4(x)}{4}  -2 \lrp{\frac{\tan^2(x)}{2} } + \ln \lrabs{\sec(x)} + C
\]
This can be rewritten as a polynomial in terms of $\sec^2(x)$ if desired. (How?)
}


\problembox{Consider the following integrals:
\begin{enumerate}[label=\alph*.]
\item $\displaystyle \int t \sin^2 t dt$
\item $\displaystyle \int t \cos^5 (t^2) dt$
\item $\displaystyle \int \tan^2 \theta \sec^4 \theta d \theta$
\item $\displaystyle \int \tan^2 (x) \cos^3 (x)dx$
\end{enumerate}
}

\ans{
\begin{enumerate}[label=$\alph*.$]
\item Use the power reduction formula for $\sin$ from \textbf{Case 3} in \textbf{Author's Note 1}, and then integration by parts.
\item Making a $u$ substitution for $u = t^2$ clears the factor of $t$, then this is a \textbf{Case 1} integral.
\item This is a \textbf{Case 1} integral from \textbf{Author's Note 2}.
\item This integrand simplifies to $\sin^2 \cos(x) $, which can be done with a $u$ substitution for $u = \sin(x)$. 
\end{enumerate}
}

\problembox{[Some Twists, or Just Good Practice] Consider the integral:
\begin{enumerate}[label=\alph*.]
\item $\displaystyle \int \frac {dx}{\cos(x) - 1}$
\item $\displaystyle \int x \tan^2 (x) dx$
\end{enumerate}
}

\ans{
\begin{enumerate}[label=\alph*.]
\item As a hint, multiply the numerator and denominator by $\cos(x) + 1$. 
\item Use $\tan^2 = \sec^2 - 1$ and then integration by parts. 
\end{enumerate}
}

\end{document}
()